\SourceLeaf{408}%
\[
\text{(2)}\quad .\ .\ .\ .\ .\quad\cos b=\cos a.\cos c+\sin a.\sin c.\cos B
\]
dividirt man (A) mit \(\sin c\), und multiplicirt (2) mit \(\operatorname{cotg}c\), und addirt, so erhält man
\[
\frac{\cos a}{\sin c}=\frac{\cos a\cos c^2}{\sin c}+\sin b.\cos A+\sin a.\cos c.\cos B
\]
also
\[
\text{(3)}\quad .\ .\ .\quad\cos a.\sin c=\sin b.\cos A+\sin a.\cos c.\cos B
\]
ebenso
\[
\text{(4)}\quad .\ .\ .\quad\cos c.\sin a=\sin b.\cos C+\sin c.\cos a.\cos B
\]
\par
Multiplicirt man (4) mit \(\cos B\), und addirt (3)
\[
\text{(5)}\quad .\ .\quad\cos a.\sin c.\sin B^2=\sin b.\cos A+\sin b.\cos B.\cos C
\]
ebenso
\[
\text{(6)}\quad .\ .\quad\cos a.\sin b.\sin C^2=\sin c.\cos A+\sin c.\cos B.\cos C
\]
\par
Multiplicirt man (5) mit \(\frac{\sin c}{\cos a}\) und zieht davon (6) mit \(\frac{\sin b}{\cos a}\) multiplicirt ab, so erhalten wir
\[
\text{(7)}\quad .\ .\ .\ .\ .\quad\sin c^2.\sin B^2-\sin b^2.\sin C^2=0
\]
oder
\[
\text{(B)}\quad .\ .\ .\ .\ .\ .\quad\sin c.\sin B=\sin b.\sin C
\]
und wenn man (A,\,5) mit \(\sin b\) dividirt, und davon (B) mit \(\frac{\cos a\sin B}{\sin b}\) multiplicirt abzieht
\[
\text{(C)}\quad .\ .\ .\ .\quad\cos A=-\cos B.\cos C+\sin B.\sin C.\cos a
\]
endlich wenn man (A,\,4) mit \(\sin c\) dividirt, und (B) mit \(\frac{\operatorname{cotg}C}{\sin c}\) multiplicirt davon abzieht
\[
\text{(D)}\quad .\ .\ .\ .\quad\operatorname{cotg}c.\sin a-\operatorname{cotg}C.\sin B=\cos a.\cos B
\]
\par
Aus diesen vier Grundformeln folgen die sogenannten \textsc{Neper}schen Analogien, und die Abkürzungen, welche durch die Bedingung, dass das sphärische Dreieck rechtwinklig sein soll, angebracht werden können, von selbst.

\{Man vergleiche mit dieser Entwickelung die von \textsc{Lagrange} im sechsten Heft des \textit{Journal de l’Ecole Polytechnique}.\}
